% Chapter 10 -- assistant strip and MCP server. Update with chat/*, gui/chat_panel.py,
% gui/settings_pages.py (AssistantPage), mcp_server.py (TOOLS), gui/mcp_dialog.py,
% config.py (mcp_enabled, assistant), suggest/runner.py (the tool-less start) and .mcp.json.example.
\chapter{Assistant and MCP}\label{ch:mcp}

Two features let an AI assistant work with your ReverseHire data, and both run through the
same door: the \emph{MCP server}. It is \textbf{off by default}, and nothing in
ReverseHire needs it.
\begin{itemize}
  \item The \textbf{assistant strip} at the bottom of the window: you ask a question and
        your own Claude Code or Codex answers from your data.
  \item The \textbf{MCP server} for any other AI client (Claude, Codex, Gemini) that you
        configure yourself.
\end{itemize}

\section{The assistant strip}
It is always visible at the bottom of the main window: one row with the toggle that shows
the conversation, the choice of assistant (\emph{Claude Code} or \emph{Codex}), the
question field, \ui{Send}, \ui{Stop}, \ui{New conversation} and the assistant settings,
and under it a status line. Opening the conversation shows your questions, the answer as
it arrives and a small notice each time a ReverseHire tool is used. Its buttons are
also in \ui{Tools $>$ Assistant}.

\textbf{It does not contain a model.} Each question starts one process of the CLI you have
installed (\code{claude} or \code{codex}), which reads the app's data through the
ReverseHire MCP server; follow-up questions resume the same conversation. Therefore:
\begin{itemize}
  \item \textbf{MCP must be on} (the dot at the top right). When it is off the strip says
        so and how to switch it on; the assistant cannot read anything meanwhile.
  \item \textbf{The CLI must be installed and signed in.} The app only runs
        \code{-{}-version} and the CLI's own sign-in check; it never makes a model call to
        test, and never keeps, logs or shows the account name. The settings page shows
        found / version / signed in for each, with \ui{Check again}, and lets you give an
        explicit path.
  \item \textbf{Usage counts against your own plan} with that provider; a notice says so
        once.
\end{itemize}

\textbf{Locked down.} The CLI is started so that only ReverseHire's tools can run: with
Claude Code your own settings, hooks and plugins are ignored, no built-in tool is
available and anything outside the ReverseHire server is refused without asking; with
Codex your own configuration is ignored, the sandbox is read-only and shell, web and
browser features are off. The question travels on standard input, never on the command
line; the process runs in an empty folder owned by the app. The reply language is the
interface language, and the assistant is told not to invent facts and to say plainly when a
tool returns nothing.

\begin{infobox}
If the strip reports that the MCP server ``did not start'' or that the assistant is
``not signed in'', the message names the fix: install the \code{mcp} package for the
Python that runs the app, or run the CLI's login command in a terminal and sign in.
\end{infobox}

\begin{infobox}
The same CLI is started a second time, \emph{without} any ReverseHire tool, when you ask for company suggestions
(chapter~\ref{ch:interest}): then it receives only the anonymised brief, in an empty folder, and cannot read your data.
\end{infobox}

\section{The MCP server}
\code{mcp\_server.py} offers the app's functions over the Model Context Protocol, on
stdio. It needs the optional \code{mcp} package (\code{pip install mcp}, or the \code{mcp}
extra). Each tool is a plain function that returns a JSON dictionary; an expected error
becomes \code{\{"error": "..."\}}, never a traceback. Logs go to stderr, only the protocol
to stdout.

{\small
\begin{tabularx}{\textwidth}{@{}l Y@{}}
\toprule
\hdr \thead{Tool} & \thead{What it returns} \\
\midrule
\code{analyze\_cv}          & Profile of a CV file: roles, skills, languages, degrees, experience, gaps, mobility, sectors. \\
\code{geocode}             & Latitude/longitude of an address, town or coordinates. \\
\code{list\_sectors}        & Sector keys with Italian and English names. \\
\code{find\_companies}      & Companies around a place (radius 0.1--50\,km, 1--100 companies), ranked, with source statuses and attribution. \\
\code{analyze\_company}     & One website: hiring status, careers page, postings; with a CV also all scores, requirements and the gate. \\
\code{saved\_searches}      & Searches saved by the app, newest first, with their id. \\
\code{search\_results}      & Results of a saved search, hiring first then by match; optional \emph{hiring only}. \\
\code{company\_card}        & A saved card by name or URL (identity status, provenance, blanks); \code{build=True} builds it first. \\
\code{list\_company\_cards}  & All saved cards: name, website, identity status, VAT, REA, ATECO, hiring status. \\
\bottomrule
\end{tabularx}
}

\section{The on/off switch}
The dot at the top right of the window shows the state (red ``MCP off'', green ``MCP on'');
clicking it, or \ui{Tools $>$ MCP server\ldots}, opens the MCP dialog with the switch, the
last use and ready-made snippets for Claude, Codex and Gemini built from the Python that
runs the app. The same switch sits on the Settings page \ui{Assistant and MCP}. It is
\code{mcp\_enabled} in \code{config.json}; it applies from the next tool call, no restart.
While off, the server still starts and lists its tools, but every call answers with how
to turn it on. Each served call records only the time and a counter in
\code{mcp\_usage.json}, never data.

\section{Connecting another client}
Copy \code{.mcp.json.example} to your client's configuration and replace
\code{<PROJECT\_PATH>} with the repository folder (on Linux/macOS use
\code{.venv/bin/python}). The app never edits other programs' files; the strip needs no
such configuration, because it starts the server itself.
\begin{lstlisting}
{ "mcpServers": { "reversehire": {
    "command": "<PROJECT_PATH>\\.venv\\Scripts\\python.exe",
    "args": ["-m", "reversehire.mcp_server"],
    "timeout": 600000 } } }
\end{lstlisting}
The long timeout matters: the first assessment may download the scoring model and a
card build reads a site slowly.

\begin{warnbox}
With MCP on, what the tools return (profile values, companies, scores, saved searches)
reaches the cloud model you use, in the strip as well as in any client. By default the
CV text, evidence lines, file path, name, contacts and home town are \emph{not} returned;
\code{include\_personal\_data=True} adds them. Your notes, contacts and tags on cards are
never returned. Switch MCP off when you are not using it.
\end{warnbox}
